%%%%%%%%%%%%%%%%%%%%%%%%%%%%%%%%%%%%%%%%%
% ANNEXES
%%%%%%%%%%%%%%%%%%%%%%%%%%%%%%%%%%%%%%%%%

\newpage
\section*{Annex A --- Inside Out Debate Format: First Design Iteration}
\addcontentsline{toc}{section}{Annex A --- Inside Out Debate Format: First Design Iteration}

The \textbf{Inside Out} format is an experimental debate format, designed using best practices from tabletop RPG game design. It seeks to be a more \textit{fun}, \textit{performance-growth oriented}, and \textit{real-life skill-building} alternative to traditional competitive debate formats. Rather than winning arguments, debaters explore values, make meaningful decisions, and earn recognition for the skills they demonstrate.

\subsubsection{Team Structure}

Debaters are organised in four teams of two debaters each and are assigned positions within the debate at random:

\begin{itemize}
    \item \textbf{Team S (Subject):} Represents the Subject of the Motion --- the person who must face a decision. Must weigh all values and choose.
    \item \textbf{Teams A, B, and C (Value Defenders):} Each argues that a particular Value held by the Subject should be prioritised in the decision, and consequently, what the Subject should decide.
\end{itemize}

\subsubsection{The Motion}

The Motion outlines:

\begin{itemize}
    \item What is the decision that has to be made.
    \item What is the background of the Subject and other relevant practical factors that can be taken into account within the decision.
    \item What are the three Values of the Subject, each assigned to one Value Defender team.
\end{itemize}

It is important to note that the Motion does not prescribe to the Value Defenders a particular position on what decision the Subject has to make. They have to construct it themselves.

\subsubsection{Speech Order}

The four teams hold alternating speeches in the following order:

\begin{center}
A1, B1, C1, S1, A2, B2, C2, S2
\end{center}

Within their speeches, the Value Defenders argue for the relevance of their value, construct a specific decision that best honours it, and engage with the arguments forwarded by the other Value Defenders.

Within their speeches, the Subject team must outline a strong reasoning that engages with and explains the interplay between all arguments of the Value Defenders, picking a particular decision. Within S1 they pick a preliminary decision and within S2 a final decision.

\subsubsection{Evaluation}

The debate is overseen by an Evaluator, who performs two tasks:

\paragraph{Commendations.} The Evaluator awards Commendations to debaters. Commendations represent acknowledgements of useful skills, knowledges, or attitudes. A list of awardable commendations is constructed by the Chief Evaluators of the Debate Tournament, along with the total number of Commendations that Evaluators can offer within that Tournament. Example commendations include Believable Passion, Witty Solution, Argument Ingenuity, and Resilience (see Section~9 of this Annex for the full list).

Over time, debaters who gather a particular commendation at least three times earn a Badge. Gathering it additional times upgrades the badge, an unlimited number of times, but each upgrade requires gathering three more Commendations than last time.

\paragraph{Oral Feedback.} The Evaluator delivers an Oral Feedback session, in which they explain to each team the commendations they received, the reasoning behind them, and give advice on how each debater can improve in the future based on their performance.

\subsubsection{Example Motion}

\begin{quote}
You are the parent of a 16-year-old teenager. Your family is middle class and lives in the suburbia of a metropolis. You have higher education. You live in a country in which Higher Education is free, but there are only a few universities that are deemed as good. The year is 2026. You evaluate that your kid is talented in music, is a good leader and has had very good STEM skills, but is also quite distracted, does not have many friends and has failed to learn a foreign language. Your kid expressed a moderate interest in studying Musicology at university after they finish high school, but is also considering not going to higher education. You have a good relationship with your child. They asked you to suggest what they should do. What do you recommend your kid they should do after they finish high school?

\medskip

Three Beliefs that you hold, relevant for this decision, are:
\begin{itemize}
    \item In life, financial stability and independence is essential for a happy life.
    \item It is important to be a well-rounded, cultured individual.
    \item The people that you surround yourself with end up defining who you are.
\end{itemize}
\end{quote}

\subsubsection{Comparison with Traditional Formats}

The table below contrasts Inside Out with World Schools debate along six dimensions.

\begin{center}
\begin{tabular}{|l|p{5cm}|p{5cm}|}
\hline
\textbf{} & \textbf{World Schools} & \textbf{Inside Out} \\
\hline
Teams & 2 teams $\times$ 3 members & \textbf{4 teams $\times$ 2 members} \\
\hline
Sides & Proposition vs.\ Opposition & \textbf{3 Value Defenders + 1 Subject} \\
\hline
Motion style & ``This House\ldots'' policy/value motions & \textbf{Personal dilemma with background \& values} \\
\hline
Goal & Win the argument --- defeat the other side & \textbf{Explore values --- help the Subject decide} \\
\hline
Evaluation & Scored on Content, Style \& Strategy (points) & \textbf{Commendations for skills \& attitudes (badges)} \\
\hline
Outcome & Win / Loss + speaker rankings & \textbf{No win/loss --- growth-oriented feedback} \\
\hline
\end{tabular}
\end{center}

\medskip

\noindent\textbf{Key shift:} Inside Out moves from adversarial persuasion to collaborative deliberation. There is no winner --- instead, debaters earn commendations that recognise individual skills, and the Subject synthesises all perspectives into a reasoned decision.

\subsubsection{Player Instructions}

\paragraph{Value Defenders.}
\begin{itemize}
    \item Formulate arguments for why your value is \textbf{important, useful, and relevant} to this dilemma.
    \item Propose a \textbf{specific decision} that best honours your value --- argue why it is right and leads to good outcomes, especially through the lens of your value.
    \item Engage with other Defenders' \textbf{values}: show yours is more important, more relevant to the dilemma, or safer to follow.
    \item Engage with other Defenders' \textbf{decisions}: attack, support (claiming it maximises your value), or endorse with addendums --- choose your positioning wisely.
    \item Your core mission: \textbf{define what honouring your value looks like} in practice, and persuade the Subject it is the right path.
\end{itemize}

\paragraph{The Subject.}
\begin{itemize}
    \item \textbf{Enter without an opinion.} You are completely torn --- you are here to listen and be persuaded, not to argue.
    \item Listen as if you \textit{are} the Subject: \textbf{let the speeches move you} as they would move this person.
    \item Also take a mental step back: \textbf{evaluate how persuasive each argument would be} for the Subject, and adjust your lean accordingly.
    \item Take notes on arguments and their interplay --- record what the Defenders say, \textbf{not your own personal thoughts}.
    \item \textbf{Collaborate quietly} with your partner: you are both embodying the same person, not competing.
    \item In your speeches, share which decision you are leaning toward and \textbf{why}, weighing all arguments to the extent they were persuasive. You must pick a specific decision, not merely name a value.
\end{itemize}

\paragraph{The Evaluator.}
\begin{itemize}
    \item Analyse the debate and each player's performance. Decide which \textbf{commendations} to award --- aim for one per debater (the same commendation can go to more than one person).
    \item Approach the round as an \textbf{educator} helping each player grow, not as a judge declaring a winner.
    \item Open with a \textbf{narrative of the debate} from your outside perspective: what each team did and how it all unfolded.
    \item Engage directly with each team: ask what they \textbf{felt they did well} and what they found difficult. Share your view and add your own observations.
    \item Ask each team what they want to \textbf{do better next time}, and give concrete advice on how to get there.
    \item Close by \textbf{awarding commendations} to each player.
\end{itemize}

\subsubsection{Rituals}

\paragraph{Rolling In.}
\begin{enumerate}
    \item \textbf{Team commitment chant} --- Each team member states their commitment aloud before the debate begins:

    Subject: \textit{``I am [character name]. I am open to being moved.''}

    Each Defender: \textit{``I believe in [value]. I will argue from that belief.''}
\end{enumerate}

\paragraph{Rolling Out.}
\begin{enumerate}
    \item \textbf{Subject's thanks} --- The Subject team thanks each Defender individually: \textit{``Thank you for speaking your truth.''}
    \item \textbf{Defenders' thanks} --- Each Defender team thanks the Subject: \textit{``Thank you for listening and carrying the decision.''}
    \item \textbf{De-roling} --- Everyone stands. Each person states their real name and shakes hands with every participant: \textit{``I am [name]. Thank you for playing with me.''}
    \item \textbf{Bleed check} --- The Evaluator asks quietly: \textit{``Did anything in that debate feel close to home for anyone?''} --- a brief, optional moment to name any emotional residue before moving into feedback. No one is required to answer.
\end{enumerate}

\subsubsection{Full Commendation List}

Commendations are organised across five categories. Each commendation has three award tiers: \textit{Earned} (threshold quality), \textit{Strong} (clear excellence), and \textit{Exceptional} (reserved for outstanding, rare performances). Some commendations carry role restrictions (Defenders only, or Subject only) as noted.

\medskip
\noindent\textbf{Style} \textit{(6 commendations --- delivery, presence, and rhetorical quality)}

\paragraph{Believable Passion.} You were truly invested in your point, showing genuine, believable conviction and care for the value you defended. Distinct from Honest Connection (self-disclosure) and Broad Range (technical delivery control).
\begin{description}
    \item[Earned] Tone conveyed genuine care at key moments; voice quality shifted naturally, eye contact held at key claims, word choices felt spontaneous.
    \item[Strong] Conviction was felt in the substance of the speech; the room or the Subject visibly reacted.
    \item[Exceptional] The delivery made the dilemma feel real and worth deeply caring about. Reserve for speeches where the affective dimension was genuinely transformative.
\end{description}

\paragraph{Entertaining Discourse.} The speech was genuinely enjoyable through wit, lightness, or a well-timed observation that made the room smile or lean in. Distinct from Poetic Appeal (literary craft).
\begin{description}
    \item[Earned] At least one moment was distinctly enjoyable: a clever comparison, a light touch, or a well-timed observation.
    \item[Strong] Liveliness was woven into the argumentation consistently; humor or vividness served to make points land more effectively.
    \item[Exceptional] Entertainment and argument were seamlessly fused throughout, reaching the level of a sharp public talk where wit and substance were inseparable.
\end{description}

\paragraph{Broad Range.} Your delivery showed impressive stylistic range, varying tone, rhythm, tempo, and volume to serve the speech. Distinct from Believable Passion (authenticity of feeling); rewards technical delivery control.
\begin{description}
    \item[Earned] At least two noticeably different delivery registers were used; the speech did not feel flat.
    \item[Strong] Variation was deliberate and purposeful: shifts in pace or tone coincided with structural moments in the argument.
    \item[Exceptional] The full palette was deployed with mastery; the debater controlled the room's attention through delivery alone.
\end{description}

\paragraph{Honest Connection.} Something genuine about yourself entered the speech --- a real disclosure rather than a rhetorical move --- that created contact between the debater and the room. Distinct from Believable Passion (conviction); rewards self-disclosure and personal presence.
\begin{description}
    \item[Earned] At least one moment of genuine self-disclosure (an admission, first-person reflection, or personal statement of feeling) that felt real rather than strategic.
    \item[Strong] The personal register appeared at key structural moments; eye contact moved around the room, direct second-person address, or deliberate pause.
    \item[Exceptional] The speech was uniquely theirs; you felt the debater was genuinely present in the room with you.
\end{description}

\paragraph{Poetic Appeal.} Language was used with aesthetic care, where rhythm, imagery, or metaphor elevated the speech as a piece of writing, independent of its entertainment value. Distinct from Entertaining Discourse (enjoyment response); rewards literary and rhetorical craft.
\begin{description}
    \item[Earned] At least one passage was noticeably well-crafted: a sentence that would stand alone as good writing, or a rhythmic construction that made a point feel inevitable.
    \item[Strong] Poetic quality sustained across multiple moments; rhetoric served the argument.
    \item[Exceptional] The speech had genuine literary quality; specific lines are still in your head afterward. Language itself became part of the argument.
\end{description}

\paragraph{True Embodiment.} You convincingly inhabited your role, arguing and deciding from inside it rather than about it from outside. For the Subject: thinking, feeling, and deciding as the character. For Defenders: arguing from genuine conviction in the value.
\begin{description}
    \item[Earned] The role's reality guided the speech throughout; stated leanings tracked the actual arguments made.
    \item[Strong] The embodiment elevated the entire round; you could see the speaker being moved, reconsidering, weighing.
    \item[Exceptional] Not applicable; omit Exceptional tier for this commendation.
\end{description}

\medskip
\noindent\textbf{Strategy} \textit{(7 commendations --- positioning, framing, and tactical intelligence)}

\paragraph{Specific Topicality.} Arguments were sharply specific to this dilemma, this Subject, and these values, with no generic claims. Distinct from Advantageous Framing (evaluative lens); rewards the content of the arguments themselves.
\begin{description}
    \item[Earned] Most substantive claims grounded in the specifics of the motion: the Subject's background, the values as stated, the actual decision at stake.
    \item[Strong] Topicality was strategic; the debater drew on details others missed.
    \item[Exceptional] Every substantive move calibrated to this specific situation; the arguments would fall apart if the Subject or dilemma were changed.
\end{description}

\paragraph{Advantageous Framing.} You defined the terms of discussion in a way that benefited your position before others could. Distinct from Specific Topicality (argument content) and Helpful Guiding (neutral Subject navigation).
\begin{description}
    \item[Earned] Explicitly established a coherent evaluation framework relevant to their value's claim.
    \item[Strong] The framing did real structural work; other debaters had to address it, and the Subject engaged with it directly.
    \item[Exceptional] The framing was original and visibly influenced the course of the debate. Award rarely.
\end{description}

\paragraph{Witty Solution.} You proposed a concrete, value-informed solution to the dilemma, specific enough that it could only have come from someone defending your value. Distinct from Advantageous Framing (criteria); rewards the solution itself.
\begin{description}
    \item[Earned] The proposed solution was specific and thought through, demonstrably flowing from the defended value.
    \item[Strong] Notably smart: navigated a genuine tension between values, or revealed an angle others missed.
    \item[Exceptional] Elegant and surprising, more persuasive than alternatives precisely because of how the value was applied.
\end{description}

\paragraph{Brave Pivot.} You changed position or conceded a point honestly, and it made your case stronger.
\begin{description}
    \item[Earned] Explicitly acknowledged an opponent's point had merit, or adjusted position in response to what was said.
    \item[Strong] The pivot was strategic: conceding less important ground to better defend more important ground.
    \item[Exceptional] Changed their mind in real time under pressure, arguing why the pivot extends rather than contradicts the previous position. Award with corresponding rarity.
\end{description}

\paragraph{Efficient Alliance.} \textit{(Defenders only)} You recognized when another Defender's arguments complemented yours and built on that alignment.
\begin{description}
    \item[Earned] Acknowledged at least one point of genuine agreement with another Defender that strengthens the case for a decision.
    \item[Strong] Explicitly constructed an alliance: ``both Team A and I agree that X, which shows the Subject that the weight of argument points toward Y.''
\end{description}

\paragraph{Careful Positioning.} \textit{(Defenders only)} Your stance toward other Defenders' decisions was deliberate and chosen rather than reactive, with positive effect.
\begin{description}
    \item[Earned] Response to other Defenders' decisions was evidently chosen rather than accidental; the stance was coherently maintained.
    \item[Strong] Positioning was in service of a coherent strategy; you could reconstruct the reasoning behind every stance taken, forming a consistent picture.
\end{description}

\paragraph{Helpful Guiding.} \textit{(Subject only)} You helped the room navigate the structure of the debate clearly, accurately, and fairly. Distinct from Advantageous Framing; the Subject's role is neutral navigation, not partisan framing.
\begin{description}
    \item[Earned] At least one moment of genuine meta-level navigation, e.g.\ ``the key question before us is now X, because Y has been established and Z has been challenged.''
    \item[Strong] The guiding was structurally important, identifying a genuine turning point or consolidating arguments fairly to all sides.
\end{description}

\medskip
\noindent\textbf{Content} \textit{(9 commendations --- argumentation, analysis, and engagement with the debate)}

\paragraph{Word Economy.} You communicated clearly and efficiently; every minute of speaking time counted. Rewards form of expression, not argument substance.
\begin{description}
    \item[Earned] The speech did not waste significant time; arguments were stated without excessive preamble, redundant restatement, or filler.
    \item[Strong] Concision actively served clarity: complex ideas were expressed with precision; understood fully on first hearing.
    \item[Exceptional] The debater said something hard to say in few words; specific lines had a density of meaning that rewarded attention.
\end{description}

\paragraph{Rebuttal Ingenuity.} Your rebuttal identified a non-obvious flaw in an opponent's argument: a hidden assumption, internal tension, or unconsidered consequence. Distinct from Argument Ingenuity (novel constructive arguments) and Comprehensive Breadth (coverage).
\begin{description}
    \item[Earned] Identified a non-obvious weakness rather than the first available attack.
    \item[Strong] Engaged with the strongest version of the argument and demonstrated why it fails even there.
    \item[Exceptional] Changed the argumentative landscape; the rebuttal was difficult or impossible to satisfactorily respond to.
\end{description}

\paragraph{Rebuttal Prioritization.} You made deliberate, strategic choices about where to concentrate your rebuttal effort, and those choices were correct. Distinct from Comprehensive Breadth (coverage); rewards deliberate focus over breadth.
\begin{description}
    \item[Earned] Concentrated force on the most important clashes, explicitly or implicitly flagging the allocation.
    \item[Strong] Prioritization revealed genuine strategic intelligence; the choice of what to ignore was as deliberate as what was attacked.
\end{description}

\paragraph{Rebuttal Accuracy.} You represented opponents' arguments faithfully before responding --- no mischaracterization, no straw-manning. No Exceptional tier; accuracy is a threshold quality. Award only when faithful representation was genuinely non-trivial.
\begin{description}
    \item[Earned] Paraphrased opponents' arguments correctly; engaged with the most charitable understanding.
    \item[Strong] Engaged with specific moments in specific speeches; accuracy made the rebuttal more effective, not just fairer.
\end{description}

\paragraph{Robust Logic.} Your arguments were well-constructed: valid, plausible, and without significant logical gaps. Distinct from Impressive Argumentative Depth; logic asks whether the reasoning chain holds structurally.
\begin{description}
    \item[Earned] Main arguments had clear claim-warrant-impact structure; a reasonable person could follow without supplying missing steps.
    \item[Strong] Logic held up under challenge: no internal contradictions, no steps requiring the audience to take things on faith.
    \item[Exceptional] Genuinely tight; would require a sophisticated objection to dislodge; obvious objections anticipated within the argument itself.
\end{description}

\paragraph{Impressive Argumentative Depth.} You went beyond the surface, interrogating implications, mechanisms, or philosophical roots. Distinct from Robust Logic (structural soundness) and Argument Ingenuity (novelty).
\begin{description}
    \item[Earned] At least one argument went meaningfully beyond its initial claim, explaining why the value's demand is well-founded or what it would look like for the Subject to live by it.
    \item[Strong] Depth sustained throughout; consistently pushed past surface-level argumentation.
    \item[Exceptional] After the speech, you found yourself reconsidering something you had assumed you understood, or the argument illuminated a dimension of the value no other speech had touched.
\end{description}

\paragraph{Comprehensive Breadth.} You engaged with the full scope of the debate, including inconvenient arguments. Distinct from Rebuttal Ingenuity (sharpness) and Rebuttal Prioritization (deliberate focus); rewards attentiveness across the full debate.
\begin{description}
    \item[Earned] Addressed the most important arguments from all other Defenders, not only the weakest or most convenient.
    \item[Strong] Demonstrated careful listening throughout; the speech reflected the actual state of the debate.
    \item[Exceptional] Managed to synthetically yet explicitly cover every significant aspect of the debate.
\end{description}

\paragraph{Argument Ingenuity.} You produced an argument that was unexpected, non-obvious, and substantively illuminating. Distinct from Rebuttal Ingenuity (reactive) and Impressive Argumentative Depth (thorough development); rewards novelty and surprise.
\begin{description}
    \item[Earned] At least one argument was notably fresh: not typically heard in this kind of debate, logically coherent, and able to withstand an obvious first objection.
    \item[Strong] Ingenious in the full sense: first surprised you, then convinced you, revealing something about the value or dilemma that prior thinking had missed.
    \item[Exceptional] The argument changed the debate; others had to respond to it and it became a fixed reference point. Very rare, award accordingly.
\end{description}

\paragraph{Accurate Synthesis.} \textit{(Subject only)} You gave a fair, accurate, and well-organized account of the full state of the debate. Distinct from Comprehensive Breadth (Defenders' engagement); rewards the Subject for faithfully describing and weighing what happened.
\begin{description}
    \item[Earned] Correctly identified the main arguments from each Defender and did not misrepresent or omit them, even where unfavorable.
    \item[Strong] Showed genuine comprehension of the debate's structure: an account of how arguments interacted, where they clashed, and what the real decision points were.
    \item[Exceptional] Identified the strongest and weakest versions of each Defender's case and grounded the final decision in a genuinely comprehensive account of the round.
\end{description}

\medskip
\noindent\textbf{Fairplay} \textit{(3 commendations --- observe team dynamics \& debrief, not only speeches)}

\paragraph{Teammate Support.} You actively helped your partner formulate something difficult in a way that visibly contributed to the team's performance. Look for: visible collaboration during prep time; a second speaker building clearly on something the first struggled with. Award when the help was substantive and its impact visible. Ask in the debrief if not fully observable.

\paragraph{Teammate Solidarity.} You showed genuine support for your partner during a hard moment, without distancing yourself. Look for: body language and verbal behavior toward a partner after a difficult moment; a second speaker who picked up where a struggling first left off without implicitly throwing them under the bus. The difficult moment must have been real and the solidarity genuine. If ambiguous, ask in the debrief.

\paragraph{Strategic Problem-Solving.} You demonstrated intelligent, adaptive problem-solving in the face of strategic challenges in the round. Look for: a visible strategy adjustment that worked; a debater who recognized a weakness in their own case and addressed it explicitly; coordinated strategic adaptation across both speeches. Both identifying and solving the problem must be visible and shared between the two team members.

\medskip
\noindent\textbf{Growth Mentality} \textit{(4 commendations --- use the debrief; always name the specific moment in feedback)}

\paragraph{Applied Feedback.} You visibly tried to apply feedback from a previous round; imperfect application still earns this. Evidence may come from direct comparison with a prior round or from the debater's own account in the debrief. What matters is the genuine attempt rather than the outcome.

\paragraph{Out of Comfort Zone.} You attempted something genuinely difficult for you and made a real go of it. Rewards the attempt, not the result. Comfort zone is individual; use the debrief to identify what was genuinely challenging. Ask directly if needed.

\paragraph{Tracking.} You showed clear awareness of your own growth across debates: where you started, what you have worked on, and where you are heading. Distinct from Applied Feedback (applying specific feedback in one round); rewards the broader self-awareness of one's development arc over time. Evidence may come from the debrief: the debater can articulate what they have improved, name skills still being worked on, and describe their trajectory with accuracy.

\paragraph{Resilience.} You encountered a genuinely difficult moment and came out intact, or stronger. The difficult moment must have been real. Award for: a strong attack responded to with a composed second speech; a delivery stumble overcome mid-speech; an unexpected move handled with adaptability. Name the specific moment in feedback; this commendation must be concrete to be useful.

%%%%%%%%%%%%%%%%%%%%%%%%%%%%%%%%%%%%%%%%%

\newpage
\section*{Annex B --- Inside Out Debate Format: Second Design Iteration}
\addcontentsline{toc}{section}{Annex B --- Inside Out Debate Format: Second Design Iteration}

The second design iteration is based on observations and participant feedback gathered during the first playtest (Jean Monet debate club, Bucharest, March 2026). The changes below supersede the corresponding elements in the first iteration (Annex A).

\subsubsection{Speech Order Changes}

The interstitial S1 speech between the first and second rounds is removed. In its place, two new elements are introduced:

\paragraph{S0 --- Subject Opening Speech.} A new opening speech before any Defender speaks. The Subject team introduces the dilemma and argues why it is genuinely hard: all three values are relevant, the decision is real, and there is no obvious right answer. This speech gives the debate proper gravitas before Defenders begin.

\paragraph{The Huddle.} A 10-minute interval inserted after C1 (and before A2). The Subject team may ask direct questions to any Defender team, grant them the floor to answer, and revoke their right to speak at any point. The Huddle allows the Subject to control the conversation toward resolving argumentative tensions, forcing arguments that need to interact to interact.

\paragraph{Thinking Time.} A 2-minute silent thinking period is added before S2, giving the Subject team time to integrate all arguments before delivering their final decision.

The new speech order is:

\begin{center}
S0, A1, B1, C1, [The Huddle], A2, B2, C2, [2-min thinking time], S2
\end{center}

\textbf{Points of Information (POIs)} are introduced during all speeches: any debater may signal to the current speaker to ask a brief question.

\subsubsection{Motion Changes}

\begin{itemize}
    \item The Subject character description in the motion should include a longer list of the Subject's personal beliefs and past decisions, to flesh out the character more fully and enable value-specific argumentation.
    \item Value Defenders should address the Subject in the second person (``you, as the Subject, should\ldots'') rather than in the third person, reinforcing the personal and direct character of the deliberation.
\end{itemize}

\subsubsection{Evaluation Changes}

Commendations are awarded in private: after the debate, the Evaluator approaches each team separately rather than delivering commendations in a public session. The oral feedback narrative and growth-oriented debrief remain public.

\subsubsection{Speech Instruction Clarifications}

\begin{itemize}
    \item Subject instructions explicitly state that the Subject must pick a specific \textit{decision} (a concrete course of action), not merely name a value they find most important.
    \item More explicit guidance is provided on what the Subject team should prepare during prep time: for S0, they should prepare to explain the dilemma's genuine difficulty and why all values are relevant; for S2, they should prepare to synthesise all arguments and justify a final specific decision.
\end{itemize}

%%%%%%%%%%%%%%%%%%%%%%%%%%%%%%%%%%%%%%%%%

\newpage
\section*{Annex C --- Debrief Form}
\addcontentsline{toc}{section}{Annex C --- Debrief Form}

The following questionnaire was administered to participants after the first playtest. Questions aim to gather data on entertainment value, format distinctiveness, participant intentions, and transformative effects.

\begin{enumerate}
    \item How fun was the debate? (scale 1--10)
    \item How fun was it compared to a debate in a traditional format? (5-choice scale)
    \item Was there anything surprisingly entertaining or engaging about it?
    \item Was there anything about the format that was quite difficult, to the point of frustration?
    \item Was there anything about the rules of the format that was confusing?
    \item Would you play this format some more? (Yes / No)
    \item If you would play it some more, in what circumstance would you?
    \item If there was to exist an alternative debate circuit centred around this format, would you attend those competitions? Even at the expense of competitions in the WS or BP format?
    \item Do you believe you played differently than you would have in a debate in WS/BP with a rather similar motion? If so, how?
    \item Do you feel different than how you usually feel after a debate? If so, how?
    \item Did you perceive yourself and others speaking similarly or differently than how you do in a regular debate? If so, how?
    \item Were you to redo this debate, what would you do differently?
    \item Were you to play more in this format, but on different motions, what would you do differently?
    \item What do you believe you learned from this debate?
    \item What do you believe this debate format does better, as a learning/empowerment tool, than WS or BP? Do you believe it does something worse?
\end{enumerate}

%%%%%%%%%%%%%%%%%%%%%%%%%%%%%%%%%%%%%%%%%

\newpage
\section*{Annex D --- Playtest Motion}
\addcontentsline{toc}{section}{Annex D --- Playtest Motion}

The following motion was used in the first playtest at the Jean Monet debate club, Bucharest, March 2026.

\begin{quote}
You are Alex. You are a lawyer. You are close to your older brother Matei. He is a very caring person to everyone, always eager to help even strangers. He always helped you when you were in need. However he was always quite dishonest with you, and you caught him lying even about important things. Matei is also an alcoholic and you believe he is quite bigoted.

\medskip

Matei was involved in a hit and run recently. He hit a cyclist that ended up hospitalised. The cyclist remembered his license plate only partially and he is now being investigated by the police. From your knowledge they only have circumstantial evidence, but he still risks going to jail. He asked you to lie that he was with you at home at the time of the accident, in order to provide him an alibi. Your assessment as a lawyer is that this will almost certainly exonerate him. However, a very small risk remains that additional material evidence could be found, at which point you might become an accessory to the crime.

\medskip

\textbf{What should you do?}

\medskip

Three Values that you hold, relevant for this decision, are:
\begin{itemize}
    \item You believe that helping people in need is very important, and a moral duty.
    \item You believe in justice: that people should receive what they deserve.
    \item You believe that family should stick together and have each other's back.
\end{itemize}
\end{quote}

%%%%%%%%%%%%%%%%%%%%%%%%%%%%%%%%%%%%%%%%%

\newpage
\section*{Annex E --- Informed Consent Form}
\addcontentsline{toc}{section}{Annex E --- Informed Consent Form}

\bigskip

I, \underline{\hspace{7cm}} of \underline{\hspace{1.5cm}} years of age, voluntarily accept to participate in a playtest followed by an interview session, both moderated by Andrei Dan Olaru, as part of his \textit{``Applying Transformative Game Design Theory to Competitive Debate''} study. This study seeks to explore the creation of an alternative debate format, informed by modern transformative game design theory. I understand that:

\begin{enumerate}
    \item This session implies video recording of the entire debate playtest and the subsequent collective interview, plus some additional questions using an online questionnaire.
    \item Pseudonyms will be used when analysing the data for confidentiality, and the researchers commit to not revealing the identities of the participants. Individual raw data may be discussed in class or with the supervisor without the participant's name to trace back to. The data obtained will be analysed by grouping and some textual examples will be used anonymously.
    \item The recording will be safely stored in the university's OneDrive to maintain the saved raw interviews protected. The video will be erased one month after the dissertation's results. The transcripts will be erased up to one year after the dissertation's results.
    \item None of the questions will go against my integrity, no harm will come to me during this session, no answer is right or wrong and I may choose not to answer any question(s) without consequences of any kind.
    \item I am free to walk out of the playtest and interview at any time if I deem it necessary without prejudice or any consequences.
\end{enumerate}

For any additional information or difficulty, the participant may contact the researcher Andrei Dan Olaru at \texttt{andrei-dan.olaru.0744@student.uu.se}.

\bigskip\bigskip

\noindent
\begin{tabular}{p{6cm} p{1cm} p{6cm}}
\underline{\hspace{6cm}} & & \underline{\hspace{6cm}} \\
Participant's signature & & Researcher's signature \\
\end{tabular}

\bigskip\bigskip

I, \underline{\hspace{7cm}} have been informed of the conditions and freely accept to participate in the research by Andrei Dan Olaru.

\bigskip

\noindent
\begin{tabular}{p{6cm} p{5cm}}
\underline{\hspace{6cm}} & \underline{\hspace{3cm}} of March 2026 \\
Participant's signature & Date \\
\end{tabular}
